%======================================================================
\section{들어가며}

\begin{yoji}
토큰 주인은 어느 계약에 허락을 줄지 정해야 해요. 다른 주인들이 이미 어느 계약에 허락을 줬는지는 공개 장부에 남아 있어요. 이 원고는 허락 기록과 송금 기록 가운데 어느 쪽이 계약의 평판 점수로 쓸모 있는지를, 기간 밖의 자료와 원시 차수라는 기준선으로 따져요.
\end{yoji}

\subsection{허락 하나로 잔고 전체가 위험해질 수 있어요}

디파이는 대출, 거래, 자산 관리를 공개 블록체인 위의 스마트 계약으로 해요~\cite{schar2021defi,werner2022sok}. 이런 서비스는 대부분 허락에서 시작해요. 토큰 주인이 어떤 계약에게 자기 토큰의 일부를 옮겨도 된다고 허락하는 거예요. ERC-20 토큰은 이 허락을 허락(allowance) 값으로 저장해요. 주인은 \texttt{approve}를 부르거나 EIP-2612의 \texttt{permit}에 서명해서 이 값을 정해요~\cite{vogelsteller2015erc20,lundfall2020permit}.

허락은 다 쓰거나, 새 허락으로 바꾸거나, 거둘 때까지 계속 살아 있어요. 그리고 이더리움의 허락은 대부분 금액에 제한이 없는 무제한 허락이에요~\cite{wang2022unlimited}. 그래서 나쁜 계약이, 또는 정직했지만 나중에 뚫린 계약이, 주인이 허락한 토큰의 잔고를 통째로 가져갈 수 있어요. 피싱 사이트는 사용자가 바로 이런 허락에 서명하도록 이끌고~\cite{he2023txphishscope}, 그렇게 하는 도구 묶음은 서비스로 팔려요~\cite{he2025drainer}. 그래서 주인과 주인의 지갑 프로그램, 그리고 다른 계약 위에 서비스를 짓는 프로토콜은 모두 같은 질문을 안고 있어요. 어느 스펜더 계약이 허락을 받을 만한가요?

\subsection{지금은 계약을 어떻게 믿나요}

지금 스펜더 계약에 대한 믿음은 몇 가지 길로 생겨요.
\begin{itemize}
\item 대부분의 주인은 지갑이 허락 요청을 보여 줄 때 결정해요. 그런데 지갑이 허락을 보여 주고 확인하는 방식에 빈틈이 있어서, 바로 그 순간이 사기에 열려 있어요~\cite{adamczyk2025approval}.
\item 코드 분석은 계약을 그 함수로 판단해요. 예를 들어 러그 풀(돈을 모은 뒤 갑자기 빼 가는 사기)을 위해 쓴 함수를 찾아내요~\cite{lin2024crpwarner}.
\item 휴먼 패스포트(Human Passport) 같은 평판 도구는 자격 증명에 무게를 달아, 한 주소가 서로 다른 한 사람의 것인지 판단해요~\cite{humanpassportnddocs}. 이것은 키를 가진 사람에 관한 질문이라, 계약에는 바로 맞는 짝이 없어요.
\end{itemize}
이 가운데 어느 것도 다른 주인들이 이미 내린 결정의 공개 기록, 곧 많은 주인이 어느 계약에 허락을 줬는지를 쓰지 않아요.

\subsection{그래프 순위 매기기와 AWP}

그래프 순위 매기기는 이 기록을 읽는 한 가지 방법이에요. 페이지랭크는 방향이 있는 화살표를 따라, 감쇠 계수를 쓰는 무작위 걷기를 할 때 각 점에 머무는 정상 확률을 그 점의 점수로 줘요~\cite{page1999pagerank,langville2006pagerank}. 아이겐트러스트(EigenTrust)와 트러스트랭크(TrustRank)는 이 걷기를 P2P 네트워크의 신뢰와 웹 스팸에 맞게 고친 것이에요~\cite{kamvar2003eigentrust,gyongyi2004trustrank}.

계정 기반 장부에서는 지금까지 화살표가 늘 지급, 곧 송금이었어요. Do와 Do는 이더리움 거래에 페이지랭크와 호지랭크(HodgeRank)를 적용해 사회적 신용을 재요~\cite{do2023hodgerank}. Do, Do, Nguyen의 AWP(적응형 가중 페이지랭크)는 송금마다 무게를 달아요~\cite{do2023awp}.
\begin{itemize}
\item 무게는 금액이 클수록 커지지만 상한이 있어요.
\item 오래된 송금일수록 무게가 로지스틱 곡선을 따라 작아져요.
\item 걷기는 최근에 송금한 주소에서 다시 출발해요.
\end{itemize}
AWP의 저자들은 AWP를 각 주소가 받은 송금의 건수와 금액에 견줘 점검해요. 그런데 이것은 AWP의 그래프를 이루는 바로 그 송금으로 만든 통계예요.

\subsection{송금과 허락은 다른 일을 기록해요}

송금은 값이 옮겨졌다는 것만 보여 줘요. 보낸 쪽이 받는 쪽을 보증한다는 뜻은 아니에요. 라우터(거래를 중계하는 계약), 시장 조성자, 에어드롭 사냥꾼, 자전 거래자는 모두 송금을 대량으로 만들어요~\cite{cong2023washtrading,luo2025airdrop}.

허락은 다른 행동을 기록해요. 돈을 가진 쪽이 스펜더 하나를 골라, 정한 금액까지 옮겨도 된다고 허락해요. 같은 주인, 스펜더, 토큰의 새 허락은 옛 허락을 바꿔 쓰므로, 마지막 허락 하나가 둘 사이의 관계를 요약해요. 그리고 허락은 송금보다 드물어요.

하지만 허락이 꼭 믿음을 나타내지는 않아요. 거래소에서 거래하려면, 사용자는 그 라우터를 어떻게 생각하든 먼저 허락해야 해요. 피싱 계약도 정직한 계약처럼 허락을 모아요. 그래서 이 원고는 허락 화살표를 권한 부여(authorization)라고 부르고, 엔도스랭크라는 이름은 방법의 이름표로만 써요. 허락은 한 주체가 가진 여러 주소를 묶는 데 쓰인 적이 있어요~\cite{victor2020clustering}. 하지만 허락으로 계약의 순위를 매기는 평판 방법은 저자들이 아는 한 없어요.

\subsection{평가 설계의 두 가지 버릇}

어느 화살표가 계약의 평판 점수로 더 나은지는 얼마쯤 평가 설계의 문제예요. 이 분야의 두 가지 버릇이 이 질문에 답하기 어렵게 해요.

첫째, 점수를 자기 화살표로 만든 통계에 견줘 점검하는 버릇이에요. 이런 일치는 만든 방식 때문에 저절로 예상되는 일치예요. 그래서 이 일치는 점수가 재려는 생각(의도한 구성개념)과 앞뒤가 맞는지를 보여 줄 뿐, 점수가 그 밖에 무엇을 알려 주는지는 말해 주지 않아요~\cite{cronbach1955construct,messick1995validity}.

둘째, 걷기를 다른 걷기하고만 견주는 버릇이에요. 페이지랭크는 그 점의 수(들어오는 화살표 수)를 매끄럽게 다듬은 것이고, 그런 수는 이기기 어려워요. 새 화살표는 이미 화살표가 많은 점에 더 잘 붙어요~\cite{barabasi1999scaling,newman2001clustering}. 차수는 링크 예측의 표준 기준선이고~\cite{libennowell2007link}, 이더리움에서도 그래요~\cite{lin2022tracking}. 그리고 차수끼리의 상관이 약한 곳에서는 페이지랭크가 평균적으로 들어오는 화살표 수에 비례해요~\cite{fortunato2008indegree}. 그러니 퍼뜨려 만든 점수는 그 밑에 깔린 수를 이길 때만 더 드는 계산값을 해요.

\begin{puri}[“원시 차수를 이긴다”는 무슨 뜻인가요?]
페이지랭크는 화살표를 따라 점수를 여러 번 퍼뜨려요. 그러면 “나에게 허락을 준 주인 자신은 얼마나 높은 점수인가”까지 점수에 들어가요. 그 대신 계산이 비싸요. 화살표를 그냥 센 수(원시 차수)가 나중 일을 똑같이, 또는 더 잘 맞힌다면 퍼뜨리는 계산은 값을 하지 못해요. 그래서 이 원고는 각 페이지랭크를 그 페이지랭크가 다듬는 원시 차수와 직접 견줘요. 엔도스랭크는 받은 허락 수와, AWP는 송금 입차수와 견줘요.
\end{puri}

\subsection{이 원고가 하는 일}

이 원고는 아비트럼 원의 스마트 계약을 두고, 두 화살표 종류를 평판 점수로 견줘요. 같은 계약들에 다음 점수를 매겨요.
\begin{itemize}
\item 엔도스랭크: 마지막 허락 위의 페이지랭크예요. AWP의 화살표 무게를 쓰고, 고르게 다시 출발해요.
\item AWP.
\item 들어오는 원시 차수 두 가지.
\end{itemize}
모든 점수를 석 달짜리 라벨 기간 두 개 앞에서 동결하고, 그 뒤에 일어난 일로 각 점수를 시험해요. 점수를 받는 주소는 모두 계약 코드를 가진 주소예요. 외부 소유 계정은 주인, 보낸 쪽, 받는 쪽으로만 그래프에 들어가요.

설계는 단계마다 정했고, 원고는 단계마다 그 커밋을 적어요.
\begin{itemize}
\item 분석 계획: 이 연구가 로그를 하나도 추출하기 전에 정했어요.
\item 두 번째 창의 등록: 이 연구가 그 창의 로그를 추출하기 전에 정했어요.
\end{itemize}
이 두 문서가 정한 등록 분석은 금액을 모든 ERC-20 토큰의 원래 단위로 두고, 2023년 10월 1일부터의 모든 사건으로 점수를 만들었어요. 그 결과를 안 뒤에 수정 분석을 정했어요. 수정 분석은 금액을 널리 거래되는 다섯 토큰의 달러 값으로 재고, 모든 점수를 동결 직전 12개월의 사건으로 만들어요. 결과 절은 수정 분석을 보고해요. 등록 분석을 위해 정한 규칙은 재분석으로 수정 분석에 적용하고, 등록 분석의 결과는 그 옆에 함께 적어요.

기여는 세 가지예요.
\begin{enumerate}
\item \textbf{기간 밖에서 각 걷기를 그 걷기가 다듬는 원시 차수와 견줬어요(두 창).} 계획도 등록 문서도 이 비교를 정하지 않았고, 모두 라벨을 안 뒤에 계산했어요. 각 화살표 종류로 만든 점수는 나중에 생기는 같은 종류의 화살표를 가장 잘 줄 세웠어요. 하지만 어느 걷기도 자기 수보다 잘 줄 세우지는 못했어요. AWP는 나중의 새 송금자를 송금 입차수만큼 잘 줄 세우지 못했고, 엔도스랭크는 나중의 새 허락 쌍을 허락한 주인 수만큼 잘 줄 세우지 못했어요. 엔도스랭크는 계획에 정한 균등 재시작으로도, AWP의 활동 재시작으로도 그랬어요. 활동 재시작은 수와의 간격을 좁혔지만 메우지는 못했어요.
\item \textbf{주인이 허락을 받지 않는 스펜더에서 엔도스랭크를 정확한 식으로 썼어요.} 그 점수는 공통 상수 $c$에 (1 + 감쇠 계수 × 주인별로 나눈 받은 무게)를 곱한 값, 곧 $r_v=c\,(1+d\,\delta(v))$예요. 이 식은 그런 스펜더 모두에서 1\% 안으로 맞아요. 그래서 대부분의 스펜더에서 이 걷기는 무게를 다시 단 허락 수예요. 나머지 스펜더는 이 그래프에도 여러 단계에 걸친 퍼뜨림이 있다는 것을 보여 줘요.
\item \textbf{어느 그래프에도 들어 있지 않은 결과에서 두 화살표 종류를 견줬어요.} 그 결과는 무기한 선물 거래소 GMX V2에서 거래하는 계약이 나중에 청산을 피하는지예요. 규칙은 이 연구가 W1의 결과를 추출하기 전에 등록했어요. 그때 W0의 닫기 기록은 추출해 두었지만 청산으로 채점하지는 않았어요. 원고는 이 규칙을 재분석으로 수정 점수에 적용해요. 주 창에서 허락한 주인 수가 송금 수보다 나중의 청산 피하기를 더 잘 줄 세웠어요. 이 결과는 그냥 센 수의 수준에서 앞서는지를 따지는 등록 규칙을 이 재분석에서는 만족하지만, 등록 분석에서는 만족하지 않았어요. 두 페이지랭크의 쌍은 판가름 나지 않았고, 표본은 작아요.
\end{enumerate}

\subsection{짝 원고와 나눠 맡은 것}

같은 저자들의 짝 원고(아직 출판되지 않음)는 두 화살표 종류를 점마다 섞어서 걷는 걷기를 다뤄요. 다층 순위 매기기~\cite{halu2013multiplex,dedomenico2015ranking}의 생각을 따른 걷기예요. 짝 원고는 이 걷기의 빼 보기 실험, 계획과 등록 문서가 이 걷기를 위해 정한 판정 규칙 두 개, 그리고 허락 걷기가 시빌 공격에 얼마나 드러나는지도 다뤄요.

\begin{itemize}
\item \textbf{두 원고가 같이 쓰는 것:} 같은 추출, 코호트, 등록, 수정 분석.
\item \textbf{두 원고가 함께 보고하는 것:} 스펜더 홀드아웃에서 원시 차수, 엔도스랭크, AWP의 계수와 차수 대비, 엔도스랭크의 닫힌 식, 풀이 실행 시간.
\item \textbf{이 원고에만 있는 것:} 트레이더 결과에서 허락 쪽과 송금 쪽을 견준 비교(P, Q, 규칙 R1–R3).
\end{itemize}
두 원고 모두 수정 분석을 먼저 보고하고 등록 분석의 값을 옆에 둬요. 그리고 두 원고는 겹치는 값을 수정 분석과 등록 분석 모두에서 서로 같게 보고해요. 같은 W1 트레이더 코호트와 청산 없이 닫은 비율에서, 짝 원고는 자기 등록 규칙이 이름 붙인 대비만 보고해요. 섞은 걷기를 그 한 층 걷기나 AWP와 견준 대비예요. 이 원고에서는 두 한 층 걷기가 트레이더 비교의 등록 쌍으로만 들어와요.

\subsection{관련 연구}

\paragraph{온체인 평판과 위험.}
위의 두 페이지랭크 연구~\cite{do2023hodgerank,do2023awp} 말고도 여러 방법이 있어요. RiskProp은 거래 네트워크를 따라 위험을 퍼뜨려 이더리움 계정을 평가하고~\cite{lin2025riskprop}, 거래 그래프의 네트워크 임베딩은 피싱 주소를 찾는 데 도움을 줘요~\cite{wu2022phishers}. 블록체인 신뢰·평판 시스템을 살핀 조사들은 명시적인 평가와 송금 활동을 흔한 신호로 꼽아요~\cite{bellini2020trust,almasoud2020reputation}. 탈중앙 신용 평가는 장부 기록을 전망 이론~\cite{hassija2020lending}이나 기계 학습~\cite{kotb2023credit}과 묶어요. 법학 연구는 이런 점수가 대신하려던 모형만큼 불투명해질 수 있다고 경고해요~\cite{packin2024credit}. ERC-20 송금 네트워크를 다룬 한 연구는 외부 소유 계정과 스마트 계약이 서로 다른 통계 패턴을 따른다는 것을 찾았어요~\cite{mukhia2026erc20}. 계약을 따로 한 집단으로 점수 매길 까닭 하나예요. 이 방법들은 모두 계정이 보내고 받은 것으로 점수를 매겨요. 이 가운데 주인이 준 허락으로 계약의 순위를 매기는 것은 없어요.

\paragraph{신뢰 퍼뜨리기와 그 약점.}
HITS는 허브와 권위를 나눠요~\cite{kleinberg1999hits}. 아이겐트러스트는 지역 신뢰를 정규화해서, 어떤 참여자도 마음대로 자기 위치를 올리지 못하게 해요~\cite{kamvar2003eigentrust}. 트러스트랭크는 미리 검증한 씨앗 집합에서 바깥으로 신뢰가 흐르게 해요~\cite{gyongyi2004trustrank}. 고르게 순간 이동하는 걷기에서는 이런 보호가 약해요. 링크 농장끼리 동맹을 맺으면 목표 점의 페이지랭크가 올라가고~\cite{gyongyi2005alliances}, 페이지랭크처럼 대칭적인 평판 함수는 시빌 공격에 안전할 수 없어요~\cite{cheng2005sybilproof}. 두 결과는 허락 그래프에도 들어맞아요. 허락 하나에는 거래 수수료만 들고 토큰이 묶이지 않으며, 계약을 배포한 사람은 주인 주소를 거의 공짜로 만들 수 있기 때문이에요~\cite{douceur2002sybil}.

\paragraph{허락.}
EIP-20이 허락을 정하고, EIP-2612가 이것을 서명 허락으로 넓혀요~\cite{vogelsteller2015erc20,lundfall2020permit}. Wang 외는 무제한 허락이 얼마나 흔하고 얼마나 큰 값을 위험에 드러내는지 쟀어요~\cite{wang2022unlimited}. 거래 기반 피싱과 드레이너(drainer) 서비스 연구는 공격자가 허락을 얻는 방법을 보여 줘요~\cite{he2023txphishscope,he2025drainer}. Adamczyk은 허락을 이용한 사기 하나와 그 사기가 파고드는 지갑의 약점을 해부해요~\cite{adamczyk2025approval}. Victor의 자기 권한 부여 휴리스틱은 허락을 한 주체의 주소들 사이의 연결로 읽어요~\cite{victor2020clustering}.

\paragraph{기준선과 평가.}
차수는 새 링크를 예측할 때 자연스러운 기준이에요~\cite{libennowell2007link,lin2022tracking}. 그리고 페이지랭크는 흔히 들어오는 화살표 수의 단조 함수처럼 움직여요~\cite{fortunato2008indegree}. 구성 타당도는 점수가 재는 것과, 점수에 이미 들어 있는 점검과의 일치를 구별해요~\cite{cronbach1955construct}. 자료를 살피기 전에 분석을 적어 두면, 분석을 입맛대로 바꿀 여지가 줄어요~\cite{nosek2018preregistration}. 저자들은 할 수 있는 곳에서는 각 규칙을, 그 규칙으로 판정할 자료를 이 연구가 추출하기 전에 정했어요. 그리고 비교마다 그것을 언제 정했는지 밝혀요.
